\section*{Supplementary Information}

\setcounter{figure}{0}
\renewcommand{\thefigure}{S\arabic{figure}}
\setcounter{table}{0}
\renewcommand{\thetable}{S\arabic{table}}

\noindent\textit{Note on file references.}\hspace{1ex}%
Throughout this Supplementary Information, paths in
\texttt{typewriter} font (e.g.\ \texttt{tests/...}, \texttt{docs/...},
\texttt{vignettes/...}, \texttt{src/...}) refer to files in the
GraphGWAS source repository at
\url{https://github.com/jfmao/GraphGWAS}; the same paths are
preserved in the Zenodo software-release archive cited under Code
Availability. Per-experiment regenerator commands quoted in
captions and notes assume the reader has cloned that repository.

% ===================================================================
% SUPPLEMENTARY TABLES (moved from main text to stay under NG's 8-item cap)
% ===================================================================

\subsection*{Supplementary Table S1 --- Multi-omics graph state after loading on 1000 Genomes Phase~3}

\begin{table}[htbp]
\centering
\small
\begin{tabular}{lr}
\toprule
Node~/~edge & Count \\
\midrule
Variant & 70{,}691{,}875 \\
Sample & 3{,}202 \\
Gene (GENCODE v47) & 20{,}092 \\
GTEx~v8 tissue eQTL & 43{,}170{,}154 \\
\texttt{HAS\_CONSEQUENCE} edges & 38{,}940{,}085 \\
STRING \texttt{INTERACTS\_WITH} ($\geq$~700) & 230{,}850 \\
RegulatoryElement (ENCODE cCRE) & 370{,}000 \\
\bottomrule
\end{tabular}
\caption{Counts in the multi-omics knowledge graph loaded on 1000
Genomes Phase~3. Variant / Sample / Gene / Pathway / GOTerm /
RegulatoryElement are first-class typed nodes; HAS\_CONSEQUENCE,
eQTL, INTERACTS\_WITH, IN\_PATHWAY, HAS\_GO\_TERM,
IN\_REGULATORY are first-class typed edges. Same counts visualised
inline in Figure~1 of the main text.}
\label{tab:sup_graph}
\end{table}

\clearpage

\subsection*{Supplementary Table S2 --- Head-to-head benchmark: HBP, GAFM and five Bayesian baselines on 1000~Genomes chromosome~22 simulations}

\begin{table}[htbp]
\centering
\footnotesize
\setlength{\tabcolsep}{5pt}
\begin{tabular}{lcccrrr}
\toprule
\multirow{2}{*}{Method} & \multicolumn{3}{c}{Rank-\#1 rate} & \multirow{2}{*}{Mean rank} & \multirow{2}{*}{Mean PIP} & \multirow{2}{*}{Median runtime} \\
\cmidrule(lr){2-4}
                       & Strong & Weak & Functional &           &           & \\
\midrule
SuSiE                              & 21/30 [52--83\%] & 17/30 [39--72\%] & 12/30 [24--57\%] & 5.74 & 0.52 & 1.87\,s \\
FINEMAP                            & 22/30 [55--85\%] & 17/30 [39--72\%] & 12/30 [24--57\%] & 6.20 & 0.51 & 2.28\,s \\
SuSiE-inf                          & \textbf{23/30 [59--88\%]} & 17/30 [39--72\%] & 11/30 [21--54\%] & 8.00 & 0.51 & 0.26\,s \\
FINEMAP-inf                        & \textbf{23/30 [59--88\%]} & 17/30 [39--72\%] & 12/30 [24--57\%] & \textbf{5.17} & 0.51 & 0.50\,s \\
\textbf{GAFM (ours)}               & 22/30 [55--85\%] & \textbf{18/30 [42--75\%]} & 12/30 [24--57\%] & 6.83 & 0.49 & \textbf{0.08\,s} \\
\textbf{HBP (ours)}                & 21/30 [52--83\%] & \textbf{18/30 [42--75\%]} & 11/30 [21--54\%] & 6.91 & 0.34 & \textbf{0.09\,s} \\
\bottomrule
\end{tabular}
\caption{Head-to-head fine-mapping benchmark. Six methods
evaluated on the same 90 simulated 1000~Genomes chromosome~22 loci
(30 reps per scenario, identical seeds): strong
($\beta{=}0.5$, $h^2{=}0.10$, random causal); weak ($\beta{=}0.2$,
$h^2{=}0.02$, random); functional ($\beta{=}0.2$, $h^2{=}0.02$,
tissue-specific eQTL causal). Rank-\#1 cells are
hits/total with the Wilson 95\% confidence interval. Mean rank,
mean PIP, and median runtime are pooled across the three scenarios.
At $n_{\text{rep}}=30$ per cell, three-percentage-point pairwise
differences (e.g.\ HBP 70\%, GAFM 73\%, SuSiE-inf 77\%) lie within
overlapping CIs and should be read as ``no significant difference
at $\alpha=0.05$'' rather than as point-equivalence; a
100-replicate replication is forthcoming. Polyfun-proxy (SuSiE with the same eQTL
prior used by HBP/GAFM via \texttt{prior\_weights}) was run on a
separate 20-replicate weak-signal benchmark and is reported in
Figure~4d rather than co-pooled here. SBayesRC operates at
genome-wide scale ($\sim$1.2~M HapMap~3 SNPs, $N \gtrsim 10^5$) and
is not directly comparable on per-locus benchmarks. Graphical
summary in main-text Figure~\ref{fig:baselines}.}
\label{tab:sup_baselines}
\end{table}

\clearpage

\subsection*{Supplementary Table S3 --- Method portfolio: recommended GraphGWAS method by scenario}

\begin{table}[htbp]
\centering
\small
\begin{tabular}{p{0.47\textwidth}ll}
\toprule
Scenario & Recommended & Primary reason \\
\midrule
Strong signal, no annotations & SuSiE~/~FINEMAP & Well-calibrated Bayesian model \\
Strong signal~+~eQTL annotations & GAFM~/~HBP & 20--30$\times$ faster, graph output \\
\textbf{Weak signal~+~informative annotations} & \textbf{GAFM} & \textbf{27--2 wins over SuSiE} \\
Dense LD~+~annotations & GAFM & Graph prior breaks LD ties \\
Multi-locus~/~shared pathway evidence & CLGF$^{\dagger}$ & EM borrows across loci \\
Genome-wide, many traits & SBayesRC & Polygenic mixture \\
Many loci, speed-critical & HBP~/~GAFM & 0.02--0.08~s per locus \\
Biobank-scale sumstats & HBP from sumstats & Pan-UKB~/~UK Biobank-ready \\
Epistasis discovery$^{\dagger}$ & LPCE$^{\dagger}$ & 42{,}000$\times$ vs exhaustive \\
\bottomrule
\end{tabular}
\caption{Decision table mapping fine-mapping scenarios to the
recommended GraphGWAS method. All listed methods share a common
input interface (graph database, BGEN file, or Pan-UKB-style summary
statistics) and a common output (\texttt{AssociationResult} and
\texttt{CredibleSet} graph nodes). Methods marked $^{\dagger}$
(CLGF and LPCE) are implemented and theoretically supported in the
released codebase but are \emph{not} benchmarked at the depth of
HBP and GAFM in the present paper; their rows reflect their intended
target scenario, not a validated head-to-head performance claim.
GLEM (graph-latent-embedding fine-mapping) is similarly implemented
but not benchmarked here. Graphical summary in main-text
Figure~\ref{fig:selection}.}
\label{tab:sup_portfolio}
\end{table}

\clearpage

% ===================================================================
% SUPPLEMENTARY FIGURES
% ===================================================================

\subsection*{Supplementary Figure S1 --- LPCE LD-pruned epistasis (preview, under development)}

\begin{figure}[htbp]
\centering
\includegraphics[width=0.95\textwidth]{fig5_m1_epistasis.pdf}
\caption{\textbf{LPCE LD-pruned co-occurrence epistasis ---
preview only, under development.}
LPCE is implemented and theoretically supported but is not
benchmarked at the depth of HBP and GAFM in this paper (CLGF and
GLEM are likewise implemented but unbenchmarked here); the panels
below report a
search-space-reduction guarantee and a 5-replicate proof of
concept, with full benchmarking against BOOST, MDR and other
epistasis-discovery tools deferred to a companion manuscript in
preparation. Readers should treat this figure as a forward-looking
sketch, not as a validated method recommendation.
(\textbf{a})~Search-space reduction on 1000 Genomes chromosome~22
common variants at LD threshold $\tau{=}0.5$: $k(\tau) \approx 0.003$
yields a 42{,}000$\times$ reduction
($5.2 \times 10^9 \rightarrow \approx 31{,}000$ pairs) without loss
of ground-truth detection. (\textbf{b})~Interaction test p-value
distribution on 5 pure-interaction simulations. (\textbf{c})~Ground-
truth rank placement: LPCE places the interacting pair at rank~\#1 in
5/5 replicates. (\textbf{d})~Runtime per locus. See Theorem~1
(Supplementary Note~S1) for the theoretical bound. The full
epistasis benchmarking programme (multi-scenario, null FPR,
BOOST/MDR baselines, higher-order $k \geq 3$) is the subject of a
companion manuscript in preparation.}
\label{fig:sup_epistasis}
\end{figure}

\clearpage

\subsection*{Supplementary Figure S2 --- Pan-UKB cross-ancestry fine-mapping (extended panels)}

\begin{figure}[htbp]
\centering
\includegraphics[width=0.95\textwidth]{figS2_panukb_cross_ancestry.pdf}
\caption{\textbf{Pan-UKB cross-ancestry fine-mapping: extended view.}
HBP sumstats-only fine-mapping on four canonical multi-ancestry
loci (FTO~/~BMI, APOA5~/~triglycerides, LDLR~/~LDL-C, HMGA2~/~height)
$\times$ four ancestries (EUR $N{=}420{,}531$, CSA $8{,}876$, AFR
$6{,}636$, EAS $2{,}709$). Pan-UKB summary statistics were streamed
via \texttt{tabix} over HTTPS ($\sim$3~s per locus) and lifted
GRCh37~$\rightarrow$~GRCh38 for alignment with our 1KG reference
panel; ancestry-matched 1KG subsets provided the LD reference.
(\textbf{a})~Top-variant HBP PIP heatmap across locus $\times$
ancestry cells. (\textbf{b})~95\% credible-set size heatmap
($\log_{10}$-scaled). (\textbf{c})~Credible-set size vs Pan-UKB $N$
(log-log) --- the same scaling relationship demonstrated in
Supplementary Figure~S5 on 1KG simulations, here replicated on
real-biobank data.
(\textbf{d})~Top-variant PIP by locus $\times$ ancestry. Notable
results: EUR BMI/LDL/Height resolve to PIP~=~1.000 with
credible-set size~1; Height/HMGA2 is cross-ancestry robust (EUR
PIP~=~1.0, CSA PIP~=~1.0, EAS PIP~=~0.97); triglycerides/APOA5
shows two-ancestry convergence (EUR and CSA both top-rank the same
indel 11:116767975:A:AAAT).}
\label{fig:sup_panukb}
\end{figure}

\clearpage

\subsection*{Supplementary Figure S3 --- GraphGWAS CLI command hierarchy}

\begin{figure}[htbp]
\centering
\includegraphics[width=0.95\textwidth]{figS3_cli_tree.pdf}
\caption{\textbf{GraphGWAS CLI command hierarchy.} The
\texttt{graphgwas} command-line interface organises 53~commands
across 15~functional groups; for legibility the visualisation shows
44 representative commands curated into 10 thematic domains and
19 function categories. Colour coding indicates
benchmark tier: red = rigorously benchmarked in this paper
(fine-mapping); orange = previewed here, full benchmark in
companion paper (epistasis); purple = platform methods implemented
but awaiting dedicated benchmarks (heritability, multivariate,
polygenic risk scores, Mendelian randomisation, gene--environment
interaction, GNN, agent); grey = shared infrastructure
(QC, plotting, I/O, REST~/~MCP servers). All commands share the
same input interface (graph database, BGEN file, or Pan-UKB-style
summary statistics) and write outputs as typed graph nodes
(\texttt{AssociationResult}, \texttt{CredibleSet}) queryable
through the database's native query language alongside the
biological annotations that produced them. Full per-command
documentation at \texttt{docs/manual/index.md}; end-to-end worked
examples in \texttt{vignettes/}.}
\label{fig:sup_cli_tree}
\end{figure}

\clearpage

% ===================================================================
% SUPPLEMENTARY NOTES
% ===================================================================

\subsection*{Supplementary Note S1 --- Mathematical proofs}\label{snote:proofs}

This note summarises the five theorems that underlie GraphGWAS'
core methods. Full proofs with technical lemmas are deposited as
\texttt{docs/MATHEMATICAL\_PROOFS.md} in the accompanying code
repository.

\paragraph{Theorem 1 (LPCE search-space reduction).}
Let $V = \{v_1, \dots, v_M\}$ be a set of $M$ variants and
$G_\tau = (V, E_\tau)$ with $E_\tau = \{(i,j) : r^2_{ij} \geq \tau\}$.
Let $S_\tau \subseteq V$ be a maximal independent set of $G_\tau$
with $k(\tau) = |S_\tau|/|V|$. The number of epistasis pairs tested
within $S_\tau$ is at most
$\binom{|S_\tau|}{2} \approx k(\tau)^2 \binom{M}{2}$, giving a
$1/k(\tau)^2$ reduction against exhaustive search. On 1KG chr22
with $\tau{=}0.5$ we measure $k(\tau) \approx 0.003$, giving
$\approx$ 100{,}000$\times$ theoretical reduction; the empirical
reduction is 42{,}000$\times$ after the additional distance and
co-carrier filters.

\paragraph{Theorem 2 (HBP fixed-point convergence).}
The HBP update operator
$T(b) = \lambda b + (1-\lambda)[\alpha s + (1-\alpha) \pi(b)]$ on
the probability simplex, with $s = \text{softmax}(z)$ constant and
$\pi(b)$ a normalised non-negative linear propagation, satisfies
$\|T(b_1) - T(b_2)\|_1 \leq L \|b_1 - b_2\|_1$ with
$L = \lambda + (1-\lambda)(1-\alpha) \rho(M) < 1$ for $\lambda > 0$,
where $\rho(M)$ is the spectral radius of
$B_{vg} B_{gp} B_{gp}^\top B_{vg}^\top$. By the Banach fixed-point
theorem [\hyperref[snote_S1_ref1]{1}], $T$ has a unique fixed point
$b^*$ and $\{b^{(t)}\}$ converges geometrically at rate $L$. At defaults
($\lambda{=}0.5$, $\alpha{=}0.6$, $\rho \approx 0.5$),
$L \approx 0.6$, giving residual $\leq 0.6^5 \approx 0.078$ in 5
rounds---empirically convergent to four decimal places.

\paragraph{Theorem 3 (GAFM causal-variant ranking).}
Under linear LD decay ($r_{ic} = 1 - d_i/D$) and bounded cross-LD
between non-causal variants ($r^2_{ij} \leq r^2_{ic} r^2_{jc}$),
the GAFM deconvolved statistic satisfies
$u_c > u_i$ for all $i \neq c$. That is, the causal variant has
the highest unique-signal contribution after LD deconvolution.

\paragraph{Theorem 4 (Null PIP bound).}
Under the null ($\beta{=}0$), the GAFM/HBP softmax PIP satisfies
$\mathbb{E}[\max_i \pi_i] \leq e^{\sqrt{2 \log n} - 1/2} / n$. For
$n = 500$ variants this gives $\max \pi_i \leq 0.04$; for $n = 5{,}000$
the bound is $\leq 0.005$. Empirical null mean
$\max$-PIP on 100 replicates: GAFM~0.004, HBP~0.003---both within the
theoretical bound.

\paragraph{Theorem 5 (CLGF EM convergence).}
The CLGF cross-locus graph fine-mapping iteration is an instance of
expectation-maximisation (EM) on a hierarchical Bayesian model with
pathway membership informing the variant prior. Monotone increase
of the expected complete-data log-likelihood follows from
Dempster, Laird \& Rubin [\hyperref[snote_S1_ref2]{2}]; convergence to a
local optimum in
$O(\log 1/\epsilon)$ iterations follows from the compactness of the
parameter space (PIPs on the probability simplex, $\sigma$ bounded
by total PIP mass). The own-locus subtraction
$\tilde{\sigma}_k = \sigma_k - \sum_{j \in V^{(l)}} \pi_j^{(t,l)} P^{(l)}_{jk}$
is essential to prevent pathological self-reinforcement.

\paragraph{References (Supplementary Note S1).}
\begin{enumerate}[leftmargin=2.2em,itemsep=0.2em,topsep=0.2em,label={[\arabic*]}]
\item\label{snote_S1_ref1} Banach, S. Sur les op\'erations dans les
ensembles abstraits et leur application aux \'equations
int\'egrales. \textit{Fundamenta Mathematicae} \textbf{3},
133--181 (1922).
\item\label{snote_S1_ref2} Dempster, A.~P., Laird, N.~M. \&
Rubin, D.~B. Maximum likelihood from incomplete data via the EM
algorithm. \textit{Journal of the Royal Statistical Society:
Series~B (Methodological)} \textbf{39}, 1--22 (1977).
\end{enumerate}

\clearpage

\subsection*{Supplementary Note S2 --- Implementation and architecture of GraphGWAS: methods, decision tree, benchmark status, and platform scope}\label{snote:platform}

\paragraph{Scope and positioning.}
GraphGWAS is a Python package of which fine-mapping (the subject
of this paper) is the first method class rigorously benchmarked.
The full source is at \url{https://github.com/jfmao/GraphGWAS}
under the MIT licence. This note describes the platform's other
capabilities---some previewed here (LPCE epistasis), some
implemented but awaiting dedicated benchmarks in forthcoming
manuscripts, and some included as infrastructure that makes the
fine-mapping methods reproducible and extensible. We disclose the
full scope here so that readers can place the present paper's
claims within the broader codebase. \textbf{Only the fine-mapping
method class (HBP, GAFM) is rigorously benchmarked in this
paper.}

\paragraph{Package architecture.}
The Python package (\texttt{graphgwas}, 34 modules, 78 passing
unit tests) is structured around seven functional groups:

\begin{itemize}
\item \textbf{Data access}: \texttt{db}, \texttt{schema}, \texttt{config}
      (graph-database connection, schema audit, global constants);
      \texttt{bgen\_reader} (per-locus BGEN streaming);
      \texttt{panukb} (Pan-UKB tabix-over-HTTPS sumstats,
      Hail LD-BlockMatrix slicer);
      \texttt{annotations} (GENCODE~/~GTEx~/~STRING~/~ENCODE
      loaders);
      \texttt{genotype} (gt\_packed encoding, PCA import).
\item \textbf{Single-locus association}: \texttt{assoc}
      (Chi$^2$, Fisher, logistic, Firth, linear with
      quantitative-trait correction); \texttt{qc} (HWE, MAC,
      missingness, heterozygosity); \texttt{popstruct} (GRM, PCA,
      GRAMMAR+, graph-spectral PCs, GPU-accelerated);
      \texttt{phenotype}, \texttt{results}, \texttt{plots}.
\item \textbf{Fine-mapping (this paper)}:
      \texttt{finemapping\_v2}
      (GAFM,
      GLEM graph-latent-embedding fine-mapping, HBP hierarchical belief propagation,
      CLGF cross-locus EM); sumstats-only entry paths
      \texttt{l1\_finemap\_from\_sumstats},
      \texttt{hbp\_finemap\_from\_sumstats}.
\item \textbf{Epistasis (companion manuscript in preparation)}:
      \texttt{epistasis\_v2}
      (LPCE LD-pruned co-occurrence, plus motif-filtered,
      differential-subgraph, and dark-matter-pair methods).
\item \textbf{Other quantitative-genetics methods (not benchmarked here)}:
      \texttt{heritability} (6 estimators: spectral, GRM-REML,
      conductance, random-effect, etc.);
      \texttt{multivariate} (cross-trait genetic correlation,
      G-matrix, coherence, pleiotropy);
      \texttt{met} (multi-environment trials; G$\times$E);
      \texttt{prs} (polygenic risk score: classical,
      graph-pruned, pathway-weighted);
      \texttt{mr} (Mendelian randomisation: IVW, Egger, weighted
      median);
      \texttt{mpat} (gene-level MPAT); \texttt{flow} (max-flow
      pathway scoring); \texttt{spectral} (rare-variant spectral
      similarity); \texttt{sv} (structural variants).
\item \textbf{Machine-learning and agent interfaces (not benchmarked here)}:
      \texttt{gnn} (HeteroGNN export/train/explain via PyTorch
      Geometric); \texttt{agent} (LangGraph ReAct agent for
      natural-language GWAS queries); \texttt{mcp\_server} (Model
      Context Protocol tool server).
\item \textbf{Infrastructure}: \texttt{parallel} (ProcessPoolExecutor
      with staggered graph-database connections); \texttt{api} (FastAPI
      REST); \texttt{sumstats}, \texttt{summary\_import},
      \texttt{simulate} (genotype/phenotype simulation under F1,
      S1, null, cross-ancestry scenarios);
      \texttt{cli} (\texttt{graphgwas} command-line interface,
      60+ commands organised by functional group).
\end{itemize}

\paragraph{Benchmark-status table.}
Table~\ref{tab:sup_benchmark_status} lists each method class
implemented in GraphGWAS alongside its validation status in the
present paper. This distinction---between methods rigorously
benchmarked here, methods previewed for follow-up manuscripts, and
methods implemented as infrastructure without dedicated
benchmarks---is essential for honest interpretation of the
codebase.

\begin{table}[htbp]
\centering
\small
\begin{tabular}{p{0.22\textwidth}p{0.17\textwidth}p{0.21\textwidth}p{0.28\textwidth}}
\toprule
Method class & Module & Status in this paper & Benchmark record \\
\midrule
\textbf{Fine-mapping (HBP, GAFM, GAFM-MX, HBP-MX, ENS)} & \texttt{finemapping\_v2} & \textbf{Rigorously benchmarked} (Figs~2--5,~7; Sup Tables~S7--S10; Sup Fig~S4 e--g) & 90 F1 simulations, 79 weak-signal replicates, 100 null replicates, 200 PIP-calibration replicates, Pan-UKB 4 ancestries; chr22 90~reps $\times$~3~$h^2$, calibration 50~reps $\times$~3~$h^2$, null FPR 100~H0 reps, hyperparam-sensitivity 30~reps $\times$~5~perturbations $\times$~2~$h^2$; cross-species (yeast 245 + Arabidopsis 54 + IRRI rice 72 + rice grain 41 + Pan-UKB 321~leads) \\
\midrule
Fine-mapping (CLGF cross-locus EM; GLEM graph-latent-embedding) & \texttt{finemapping\_v2} & Implemented + theoretically supported (Theorem~5 for CLGF), \textbf{not benchmarked here} & Architectural preview; full benchmarks deferred to follow-up work \\
\midrule
Epistasis (LPCE; further methods forthcoming) & \texttt{epistasis\_v2} & LPCE preview only, under development (Supp.~Fig.~S1) & 5 replicates; full benchmark in companion manuscript (in preparation) \\
\midrule
Single-locus GWAS + GRAMMAR+ & \texttt{assoc},\,\texttt{popstruct} & Used internally (yeast $\lambda_{\text{GC}}$ 0.98 reported) & Validated vs PLINK2 at $r = 1.0000$ on 82k shared variants \\
\midrule
Heritability (6 estimators) & \texttt{heritability} & Mentioned only & Spectral $h^2$ computed on 35 yeast traits; not formally benchmarked vs GCTA-GREML \\
\midrule
Multivariate (r$_G$, G-matrix, pleiotropy, coherence) & \texttt{multivariate} & Not used & Implemented; 35-trait yeast r$_G$ matrix computed; no formal benchmark \\
\midrule
Polygenic risk score (classical, pathway-weighted) & \texttt{prs} & Not used & Implemented; no formal benchmark \\
\midrule
Mendelian randomisation (IVW, Egger, weighted median) & \texttt{mr} & Not used & Implemented; no formal benchmark \\
\midrule
Multi-environment trials / G$\times$E & \texttt{met} & Not used & Implemented; no formal benchmark \\
\midrule
GNN (HeteroGNN) & \texttt{gnn} & Mentioned only (Discussion future directions) & Implemented; not trained for this paper \\
\midrule
AI agent (LangGraph ReAct) & \texttt{agent},\,\texttt{mcp\_server} & Not used & Proof-of-concept implementation; no formal benchmark \\
\bottomrule
\end{tabular}
\caption{Benchmark-status table for all method classes in
GraphGWAS. Only the fine-mapping class is rigorously benchmarked
in the present paper. Eight further method classes are implemented
in the released package but either previewed for follow-up work
(epistasis) or shipped as infrastructure without formal
benchmarks. This table is the honest record of the codebase's
scope versus the paper's claims.}
\label{tab:sup_benchmark_status}
\end{table}

\paragraph{Reproducibility and extensibility.}
Every result in this paper is regeneratable by a single command
(see Code Availability and \texttt{docs/REPRODUCIBILITY.md} in the
released repository). New methods plug into the shared
\texttt{load\_locus\_variants} input contract (graph database, BGEN,
or Pan-UKB-style sumstats) and the shared
\texttt{CredibleSet}~/~\texttt{AssociationResult} output contract
(results become graph nodes queryable by Cypher alongside
biological annotations). New annotations plug into the
\texttt{graphgwas.annotations} loader via a small typed-edge
registration schema. Distribution channels: GitHub source,
Zenodo-versioned package DOI (to be assigned on acceptance),
graph-database dumps for yeast (0.5~GB) and human with full
multi-omics (17~GB).

\paragraph{User-facing documentation.}
The package ships with three levels of documentation. (i)~A
top-level \texttt{README.md} with quick-start, key features,
three-interface table (CLI / REST / MCP), and pointers to detailed
guides. (ii)~A per-command reference at \texttt{docs/manual/index.md}
listing all 53 commands organised by 15 functional groups; each
command has its own page under \texttt{docs/manual/commands/}
following a standard template (description, usage, arguments,
options, output schema, examples). The CLI hierarchy is visualised
in Supplementary Figure~S3. (iii)~End-to-end vignettes in
\texttt{vignettes/}, including \texttt{fine-mapping-quickstart.md}
which walks from Pan-UKB summary-statistics fetch through an
ancestry-resolved credible set in $\sim$15 minutes, replicating
the FTO/BMI headline result of the main paper (Figure~\ref{fig:cross_anc}). An
installation guide at \texttt{docs/INSTALL.md} documents Python
environment setup, graph-database configuration for the optional full-graph
deployment, and Hail setup for the optional Pan-UKB in-sample LD
path.

\paragraph{Three programmatic interfaces.}
The same 53 procedures are exposed through three interfaces sharing
a common implementation: (i)~the \texttt{graphgwas} CLI (Click-based,
53~commands, Supplementary Figure~S3); (ii)~a FastAPI REST server
with 37 endpoints for web or pipeline integration
(\texttt{graphgwas serve}); and (iii)~a Model Context Protocol (MCP)
server with 16 tools for AI-agent access from any MCP-compatible
client (\texttt{graphgwas mcp}). All three interfaces call the
same underlying Python module functions; there is no duplication of
logic between interfaces.

\paragraph{Pre-registered control experiments (placebo prior, null-overlap permutation, hyperparameter sensitivity).}\label{snote:controls}%
Three pre-registered control experiments were run on the GraphGWAS
chr~22 testbed to bound the claims made in the main text.

\textbf{Null-overlap permutation.} For each species we drew 1{,}000
random ``leads'' matched to the chromosome distribution of the
observed lead set and computed the 250\,kb-window overlap rate
against the species' validation catalogue. Rice: 33.3\,\% observed
vs 0.9\,\% null mean ($p < 0.001$, $\sim$38-fold enrichment).
\emph{Arabidopsis}: 5.6\,\% observed vs 0.0\,\% null mean
($p < 0.001$; catalogue too sparse for an enrichment estimate).
Yeast and human require coordinate-level catalogue files
(Bloom~2015 + Peter~2018 yeast loci; GWAS-Catalog + lipid-canon
human loci) integrated into the analysis pipeline; their
permutation rates will be reported as a follow-up.

\textbf{Hyperparameter sensitivity (10 simulated chr~22 loci, $\beta=0.5$, $h^2=0.10$).}
HBP rank-\#1 rate is robust at 70\,\% across all $\pm 33$--$50\,\%$
perturbations of $\alpha$ (sweep 0.4--0.8), $\lambda$ (0.3--0.7),
$T$ (3--7), and $r^2_{\text{smooth}}$ (0.2--0.4); a single dip to
60\,\% at $r^2_{\text{smooth}}{=}0.35$. GAFM is highly
$\alpha$-sensitive: at $\alpha=0.5$ rank-\#1 is 10\,\%; at
$\alpha=0.9$ rank-\#1 is 70\,\%; at $\alpha=1.0$ rank-\#1 is
70\,\%. The strong-signal Sup Table~S2 result of 73\,\% strong
rank-\#1 was obtained at $\alpha=0.9$ (regime-specific default);
the 27\,--\,2 weak-signal headline used $\alpha=0.5$. We disclose
the regime-specific defaults explicitly in Methods §GAFM.

\textbf{Placebo-prior control (30 simulated chr~22 loci, $\beta=0.5$, $h^2=0.10$).}
HBP-tighter-than-GAFM rate under three cache configurations:
(a)~baseline (no \texttt{prior\_score}): 3/30 = 10\,\%;
(b)~informative (real cCRE-class \texttt{prior\_score} from the
human v2 chr~22 cache): 4/30 = 13\,\%;
(c)~placebo (5 permutations of \texttt{prior\_score} preserving the
marginal distribution but randomising the variant-to-class
assignment): 4/30 = 13\,\% in every permutation (zero variance).
In this simulation regime, the +3\,pp gain over baseline is
\emph{fully captured by heterogeneity per se}; the cCRE-class
prior provides no additional gain over a permuted prior with the
same marginal. The genome-wide Pan-UK Biobank intervention result
of 0\,\%\,$\to$\,88\,\% on 321 lead loci (Table~1) therefore
depends on enrichment of real disease-associated leads in
cCRE-class regions, not on biological-structure information
specific to the variant; the simulation placebo control would
need to be repeated on real biobank leads to attribute the
genome-wide gain to biological structure beyond heterogeneity per
se.

Full numerical outputs are deposited in \texttt{results/null\_overlap/},
\texttt{results/hyperparam\_sensitivity/}, and
\texttt{results/placebo\_prior/} alongside the deferred-experiments
report at \texttt{results/deferred\_experiments\_report.md}.

\textbf{Mixture-prior controls (50 reps $\times$ 3 $h^2$ calibration; 100 H0 null FPR;
30 reps $\times$ 2 $h^2$ $\times$ 5 perturbations sensitivity, all on chr~22).}
The mixture-prior posterior reweighting (GAFM-MX, HBP-MX, ENS) sharpens
top PIPs by 2--3$\times$ at the same rank parity as base GAFM/HBP but is
\emph{anti-conservative} at high PIP bins (empirical TDR $\approx 0.62$
in the $[0.9, 1.0]$ PIP bin vs the diagonal-target $\approx 0.95$);
under the null ($y = \varepsilon$, no causal) the rate of $\max\mathrm{PIP}\geq 0.5$
is $0/100$ for GAFM/HBP, $1/100$ for SuSiE/SuSiE-inf/FINEMAP-inf,
$6/100$ for GAFM-MX and ENS, and $10/100$ for HBP-MX; at $\geq 0.9$ all
methods are $\leq 1/100$. GAFM-MX rank-1 rate and mean PIP at causal are
unchanged ($\Delta < 0.005$) under five perturbations of the SBayesRC
mixture defaults --- $\pi$-uniform $(0.25,0.25,0.25,0.25)$,
$\pi$-concentrated $(0.001,0.001,0.001,0.005)$, $\gamma$ shifted
$10\times$ left or right --- so the operational behaviour does not
depend on tuning the mixture hyperparameters within reasonable ranges
(\texttt{benchmark\_v15\_calibration\_null.py}, \texttt{benchmark\_v15\_hyperparam\_sensitivity.py};
panels e--g of Sup Fig~S4). Mixture-prior PIPs should therefore be
interpreted as an operational ranking score (sharper credible sets,
higher rank-1 rate) rather than as posterior probabilities; users
needing calibrated PIPs should report base GAFM/HBP or SuSiE outputs
alongside.

\paragraph{LPCE algorithm details (preview, under development).}\label{snote:lpce}%
LPCE (LD-pruned co-occurrence pairwise-epistasis) is one of the
implemented-but-not-benchmarked methods in the platform; the
algorithm is described here for technical reproducibility.
Given $M$ candidate variants in a region, LPCE first constructs a
maximal independent set $S_\tau$ on the LD graph
$G_\tau = \{(i,j): r^2_{ij} \geq \tau\}$, $\tau = 0.5$ by default.
The greedy constructor sorts variants by MAF descending and admits
a variant to $S_\tau$ iff it has $r^2 < \tau$ to every
already-admitted member. Candidate pairs within $S_\tau$ are
enumerated subject to (i) physical distance
$|\text{pos}_i - \text{pos}_j| > 100$~kb and (ii) co-carrier count
$\geq 5$. For each admitted pair LPCE fits
$Y = \alpha + \beta_1 G_1 + \beta_2 G_2 + \beta_I (G_1 \cdot G_2) + X\gamma$,
with $X$ the same covariates as in the single-locus association
of the Online Methods, and tests $H_0: \beta_I = 0$ by a Wald
statistic on $\beta_I$ against $\chi^2_1$. Multiple testing uses
Benjamini--Hochberg at 5\% false discovery rate (FDR). On 1KG
chr22 common variants ($M = 102{,}467$, $\tau{=}0.5$) the greedy
pruner yields $|S_\tau| \approx 250$, and after distance and
co-carrier filters $\approx$~31~K candidate pairs remain --- a
42{,}000$\times$ reduction against the exhaustive
$5.2 \times 10^9$ pairs (Theorem~1, Supplementary Note~S1). The
search-space-reduction guarantee and a 5-replicate proof-of-concept
are reported in Supplementary Figure~S1; full benchmarking against
BOOST, MDR and other epistasis-discovery tools is the subject of a
companion manuscript in preparation. Readers should treat this
section, the LPCE row of the benchmark-status table
(Supplementary Table~S4), the LPCE branch of the decision tree
(Supplementary Figure~\ref{fig:selection}), and Supplementary
Figure~S1 as a forward-looking sketch, not as a validated method
recommendation; CLGF (cross-locus EM) and GLEM (graph-latent-embedding
fine-mapping) are likewise architectural previews in the present
paper.

\paragraph{Decision tree for method selection.}
A visual summary of the scenario--method mapping in
Supplementary Table~S3 --- included as part of the implementation
and architecture of GraphGWAS --- is provided in Supplementary
Figure~\ref{fig:selection} (rendered at the end of the
supplementary figures, after Figure~S5, to preserve the existing
numbering of Figures~S4 and~S5).

\clearpage

\subsection*{Supplementary Note S3 --- Multi-omics coverage ablation (rice)}\label{snote:ablation}

This note expands on the ablation result in Results \S``\nameref{sec:heterogeneity}''.

\paragraph{Protocol.} On each of 5 simulated grain-quality causal loci in
the 3{,}000 Rice Genomes panel (GW2, BG1, GS5, TGW6, GW8/OsSPL16; each at
$\pm$250~kb around the Ren-2023 causal variant position), we took the
per-chromosome rice multi-omics graph cache (RAP-DB variant$\to$gene via
snpEff~ANN, Ren 2023 gene$\to$pathway curated catalogue of 269 grain genes,
RicePPINet gene$\leftrightarrow$gene PPI at probability~$\geq$~0.7). At
each retention fraction $f \in \{0.1, 0.25, 0.5, 0.75, 1.0\}$ we randomly
dropped each (variant, edge) entry with probability $1-f$ and re-ran HBP
on the same sumstats-derived $z$ and LD matrix. Ten independent bootstrap
seeds were used per non-unit fraction ($f = 1$ is deterministic).

\paragraph{Per-locus rank of causal variant.} Means over 10 seeds:

\begin{center}
\begin{tabular}{lrrrrrr}
\toprule
\textbf{Locus} & \textbf{Baseline PIP} & \textbf{10\%} & \textbf{25\%} &
\textbf{50\%} & \textbf{75\%} & \textbf{100\%} \\
\midrule
GW8/OsSPL16  & 0.040   &  2.2 &  2.1 &  2.0 &  2.0 &  2.0 \\
BG1          & 1.3$\times$10$^{-3}$ & 154.3 & 58.6 & 7.5 &  4.5 &  4.0 \\
TGW6         & 2.7$\times$10$^{-3}$ &  64.7 & 34.7 & 15.2 &  6.7 &  3.0 \\
GS5          & 2.1$\times$10$^{-4}$ & 514.3 & 764.0 & 1025.6 & 1156.0 & 1139.0 \\
GW2          & 6.1$\times$10$^{-7}$ &10791 &10899 &11033 &11130 &11197 \\
\bottomrule
\end{tabular}
\end{center}

\paragraph{Per-locus credible-set size.} CS size at the same fractions:

\begin{center}
\begin{tabular}{lrrrrr}
\toprule
\textbf{Locus} & \textbf{10\%} & \textbf{25\%} & \textbf{50\%} & \textbf{75\%} & \textbf{100\%} \\
\midrule
GW8/OsSPL16  &  184 &  354 &  558 &  700 &  783 \\
BG1          & 7164 & 7164 & 7164 & 7164 & 7164 \\
TGW6         & 4017 & 4047 & 4086 & 4115 & 4157 \\
GS5          & 8577 & 8577 & 8577 & 8577 & 8577 \\
GW2          &  901 & 1155 & 1472 & 1707 & 1865 \\
\bottomrule
\end{tabular}
\end{center}

\paragraph{Seed-variance.} Rice shows rank\_std $\in$ [7, 18] across 10
seeds for borderline and unresolved loci, confirming that the random
edge-dropping produces genuinely different caches whose HBP outputs
differ. A companion human-chr22 ablation on simulated phenotypes using a
GTEx eQTL + STRING PPI cache (combined\_score~$\geq$~700, nearly uniform
1-gene / $N$-PPI per annotated variant) yielded rank\_std $=$ 0 for all
5 test loci at every fraction in both strong-signal ($\lambda=6$) and
weak-signal ($\lambda=3.5$) regimes --- no coverage sensitivity. This
distinguishes the rice cache's heterogeneity (Ren-2023 hand-curated
pathway labels vs. gene-only vs. PPI-only) as the driver of the
observed signal in rice, and implies that uniform-density multi-omics
priors do not benefit fine-mapping above a saturation density.

\paragraph{Reproducibility.} Driver scripts are
\texttt{tests/multiomics\_ablation\_test.py} (reusable scaffold),
\texttt{tests/multiomics\_ablation\_rice.py} (rice driver), and
\texttt{tests/multiomics\_ablation\_human\_weak.py} (human chr22
weak-signal companion). Combined TSV output is
\texttt{data/rice\_3k/results/ablation\_rice\_combined.tsv};
full report with classification methodology is
\texttt{data/rice\_3k/results/ablation\_multi\_species\_report.md}.

\clearpage

\subsection*{Supplementary Note S4 --- Cross-species ground-truth catalogs}\label{snote:multispecies_gt}

This note documents the literature and database sources used for the
ground-truth validation of cross-species fine-mapping (Results
\S``\nameref{sec:multispecies}'').

\paragraph{Rice (Ren \emph{et al.} 2023).}
269 grain-quality genes from the comprehensive Science Bulletin
review [\hyperref[snote_S4_ref3]{3}], with chromosome and position from snpEff
ANN of the 3kRG pseudo-canonical VCF
(\texttt{data/rice\_3k/ground\_truth/grain\_quality\_causal\_genes.tsv}
+ \texttt{gene\_position\_index.tsv}).

\paragraph{Yeast (Bloom 2015 + Peter 2018 + curation).}
35 trait-to-gene mappings drawn from [\hyperref[snote_S4_ref2]{2},\hyperref[snote_S4_ref1]{1}]
plus standard yeast genetics canon. Examples:
TUP1/SPT7/ADH1-3 (ethanol tolerance), CUP1/CUP2/CTR1/ACE1 (copper),
HSP82/HSP104/HSP12/TPS1/TPS2 (heat), GAL10/GAL80
(galactose utilisation), TOR1/PDR1/PDR3/YRR1 (caffeine, multidrug
resistance), ERG3/ERG6/ERG11/PDR5/CDR1 (fluconazole/sterol),
HOG1/SLN1/HAL3/ENA1 (osmotic), TUB1-3/BEN1 (benomyl/microtubule),
RAD52/RAD53/RNR1-3/MEC1 (HU/DNA-damage), GAL/HXT/GUT/RBK1
(carbon-source). Positions parsed from
\texttt{tests/data/yeast/SGD\_features.tab}.

\paragraph{Arabidopsis (AraGWAS bonferroni + literature).}
275 SNP--phenotype rows from AraGWAS bonferroni-significant
associations (\texttt{tests/data/arabidopsis/aragwas\_bonf\_associations.csv})
plus 24 literature-canonical entries we added for traits not in
AraGWAS but with well-known causal genes:
FLC (Sheldon 1999), FRI (Johanson 2000), CONSTANS (Putterill 1995),
GIGANTEA (Park 1999), VIN3 (Sung 2004), FT (Kobayashi 1999), AP1
(Bowman 1989), DOG1 (Bentsink 2006), HKT1 (Rus 2006). The AraGWAS file is dominated by chr-4 hits; cross-chromosome
validation relies on the literature layer.

\paragraph{Human (GWAS Catalog literature canon).}
8 chr22 loci linked in the GWAS Catalog or major lipid GWAS
to our 4 traits: APOL1/APOL2 (LDL, Klarin et al.\ 2018),
SREBF2 (LDL/Height, Willer et al.\ 2013), BCL2L13/BID (LDL),
PLA2G6 (TG, GLGC 2013), PPARA (TG, Surakka et al.\ 2015),
KCNJ4 (Height). Chromosome-22-only because that is
where our graph cache is built.

\paragraph{Validation criterion.} A fine-mapped lead is "validated"
if any catalog feature for its trait lies within 250~kb of the lead
position on the same chromosome. This is the same criterion used
across all four species. The 250~kb window matches the typical
LD-block scale of the densest panel (rice 3kRG); for human it is
about half a typical LD block (Pan-UKB chr22 leads have $\geq 1$~Mb
LD) and so likely under-counts true positives.

\paragraph{References (Supplementary Note S4).}
\begin{enumerate}[leftmargin=2.2em,itemsep=0.2em,topsep=0.2em,label={[\arabic*]}]
\item\label{snote_S4_ref1} Bloom, J.~S. \emph{et~al.} Genetic
interactions contribute less than additive effects to quantitative
trait variation in yeast. \textit{Nature Communications}
\textbf{6}, 8712 (2015).
\item\label{snote_S4_ref2} Peter, J. \emph{et~al.} Genome evolution
across 1{,}011 \textit{Saccharomyces cerevisiae} isolates.
\textit{Nature} \textbf{556}, 339--344 (2018).
\item\label{snote_S4_ref3} Ren, D., Ding, C. \& Qian, Q. Molecular
bases of rice grain size and quality for optimized productivity.
\textit{Science Bulletin} \textbf{68}, 314--350 (2023).
\end{enumerate}

\clearpage

\subsection*{Supplementary Figure S4 --- PIP calibration and null false-positive rate}

\begin{figure}[htbp]
\centering
\includegraphics[width=0.95\textwidth]{fig4_calibration_null_fpr.pdf}
\caption{\textbf{PIP calibration and null false-positive rate.}
(\textbf{a})~PIP calibration across 200 F1 simulations $\times$
4 heritability levels ($h^2 \in \{0.02, 0.05, 0.10, 0.20\}$).
Observed true-discovery rate is plotted against reported PIP in
eight bins of width 0.125. Ideal calibration lies on the diagonal
(dashed). SuSiE, FINEMAP and GAFM track the diagonal; HBP is
conservative at mid-range PIPs and caps at $\approx$0.7.
(\textbf{b},\textbf{c})~Null maximum-PIP distribution across 100
null simulations (no causal variant, Gaussian phenotype) on random
50~kb yeast loci. Mean $\max$\,PIP is 0.003 (HBP), 0.004 (GAFM), 0.013
(SuSiE), 0.029 (FINEMAP); the four base methods all report
$\max$\,PIP~$<$~0.5 in every replicate. The theoretical bound
$\mathbb{E}[\max_i \pi_i] \leq e^{\sqrt{2 \log n} - 1/2}/n$
(Supplementary Note~S1, Theorem 4) predicts $\leq$0.005 at
$n = 5000$, matching empirically to 2\% accuracy. (\textbf{d})~Null
credible-set size: under $\beta = 0$, credible sets grow to cover
$\approx$90\% of the locus --- the correct response to ``no
signal''. (\textbf{e})~Calibration of the mixture-prior methods on the
chr22 single-causal $\times$ 50 reps $\times$ 3 $h^2$ panel
($h^2\in\{0.02,0.05,0.10\}$,
\texttt{benchmark\_v15\_calibration\_null.py}). The mixture-prior
methods (GAFM-MX, HBP-MX, ENS) sharpen
top PIPs by 2--3$\times$ at the same rank parity but are
\emph{anti-conservative} at high PIP bins (empirical TDR
$\approx 0.62$ in the $[0.9, 1.0]$ bin vs the diagonal-target
$\approx 0.95$); base GAFM/HBP and SuSiE/SuSiE-inf/FINEMAP-inf
remain on or near the diagonal across the full range. Users who
need calibrated PIPs should report base GAFM/HBP or SuSiE outputs;
mixture-prior PIPs should be interpreted as an operational ranking
score (sharper credible sets, higher rank-1 rate) rather than as
posterior probabilities. (\textbf{f})~Null FPR on 100 H0
loci (chr22, $y=\varepsilon$): GAFM/HBP $0/100$, SuSiE/SuSiE-inf/FINEMAP-inf
$1/100$, GAFM-MX/ENS $6/100$, HBP-MX $10/100$ at the
$\max\mathrm{PIP}\geq 0.5$ threshold; at $\geq 0.9$ all five
methods $\leq 1/100$. The mixture step's anti-conservative tail is
a well-understood property of post-hoc multiplicative reweighting
of LD-aware PIPs (Methods, \S\,``Mixture-prior posterior reweighting and ensemble''):
the multiplicative ABF redistributes mass aggressively at large
$|z|$, and a small fraction of LD-correlated noise variants under
the null still receive a modest top PIP ($\approx 0.5$--$0.7$).
(\textbf{g})~Hyperparameter sensitivity:
GAFM-MX rank-1 rate and mean PIP at causal are unchanged
($\Delta < 0.005$) under $5\times$ perturbations of the SBayesRC
defaults --- $\pi$-uniform $(0.25, 0.25, 0.25, 0.25)$,
$\pi$-concentrated $(0.001,0.001,0.001,0.005)$,
$\gamma$ shifted $10\times$ left or right --- across $30$ reps and
two $h^2$ levels. The mixture is dominated by the largest-variance
component for $|z| \geq 5$, so reasonable perturbations of $\pi/\gamma$ do
not change the per-variant ranking
(\texttt{benchmark\_v15\_hyperparam\_sensitivity.py}).
(\textbf{h})~Confidence sharpening at the same rank parity:
mean PIP at the causal variant for the chr22 single-causal F1 panel
($h^2 \in \{0.02, 0.05, 0.10\}$, 30 reps each). The mixture-prior
methods deliver $2$--$3\times$ the confidence of base GAFM/HBP at
identical rank-1 counts, with the largest sharpening at the
weakest signal where fine-mapping needs help most.}
\label{fig:sup_calibration}
\end{figure}

\clearpage

\subsection*{Supplementary Figure S5 --- HBP fine-mapping quality scales with sample size}

\begin{figure}[htbp]
\centering
\includegraphics[width=0.95\textwidth]{fig8_power_vs_sample_size.pdf}
\caption{\textbf{HBP fine-mapping quality scales monotonically
with sample size.} 1000 Genomes Phase~3 chromosome~22 subsampled
at $N \in \{500, 1{,}000, 2{,}000, 3{,}000\}$, 30 F1 simulations
per $N$ at $h^2 = 0.10$, causal MAF 10--40\%. (\textbf{a})~Mean
causal-variant PIP scales 0.54~$\rightarrow$~0.77 across the $N$
range. (\textbf{b})~Rank-\#1 rate rises from 45\% at $N{=}500$ to
73\% at $N{=}3{,}000$. (\textbf{c})~95\% credible-set size shrinks
from median 40 at $N{=}500$ to median 8 at $N{=}3{,}000$.
(\textbf{d})~Per-replicate causal-variant rank distributions by
$N$. The linear extrapolation to $N > 10{,}000$ supports the
biobank-scale Pan-UKB results in Figure~\ref{fig:cross_anc}.}
\label{fig:sup_power}
\end{figure}

\clearpage

\subsection*{Supplementary Figure S6 --- A decision tree maps fine-mapping scenarios to the recommended GraphGWAS method}

\begin{figure}[htbp]
\centering
\includegraphics[width=0.85\textwidth]{fig7_method_selection.pdf}
\caption{\textbf{A decision tree maps fine-mapping scenarios to
the recommended GraphGWAS method.} Recommended GraphGWAS method by
signal strength, annotation informativeness, LD complexity, and
whether the workload is per-locus or genome-wide. For strong
signal without annotations, SuSiE~/~FINEMAP remain the standard.
For strong signal with eQTL annotations, HBP matches accuracy at
20--30$\times$ the speed. For weak signal with informative
annotations, GAFM wins decisively (27--2 over SuSiE; main-text
Figure~\ref{fig:weak_signal}). For dense LD with annotations,
GAFM's graph prior breaks LD ties that flat-prior methods cannot.
For genome-wide workloads, SBayesRC is appropriate; HBP/GAFM
target high-precision per-locus refinement at SBayesRC-surfaced
leads. \textbf{LPCE} (LD-pruned co-occurrence
pairwise-epistasis; algorithm details in Supplementary Note~S2,
``\nameref{snote:lpce}''), shown at the epistasis branch, is
presented as an architectural preview
only --- it is implemented and theoretically supported
(Theorem~1, Supplementary Note~S1) but \textbf{not benchmarked at
the depth of HBP and GAFM in this paper}; CLGF (cross-locus EM)
and GLEM (graph-latent-embedding fine-mapping) are likewise
implemented in the codebase but are not benchmarked here either.
Full LPCE benchmarking against epistasis-discovery baselines is
the subject of a companion manuscript in preparation, and the
LPCE branch of the decision tree should be read as a
forward-looking sketch rather than a validated recommendation. All
methods are accessible through the same input interface (graph
database, BGEN, or Pan-UKB-style sumstats) and return
\texttt{CredibleSet} nodes queryable alongside biological
annotations. This figure is referenced from Supplementary
Note~S2 (``\nameref{snote:platform}'').}
\label{fig:selection}
\end{figure}

\clearpage

\subsection*{Supplementary Figure S7 --- 3kRG grain weight + shape pass:
4-trait Manhattan, Q--Q, and nine-method recovery scorecard}

\begin{figure}[p]
\centering
\includegraphics[width=0.95\textwidth]{rice_grain_manhattan.pdf}\\[1ex]
\includegraphics[width=0.50\textwidth]{rice_grain_qq.pdf}\hfill
\includegraphics[width=0.48\textwidth]{rice_grain_recovery_scorecard.pdf}
\caption{\textbf{3kRG grain weight + shape pass --- four-trait
Manhattan, Q--Q, and nine-method recovery scorecard.}
(\textbf{Top}) Genome-wide $-\log_{10}\!P$ Manhattan tracks for
TGW, GL, GW and RLW. (\textbf{Bottom-left}) Per-trait quantile--quantile
plots ($\lambda_{GC} \in [1.41, 1.78]$) with the
$y = x$ null and slope-$\sqrt{\lambda_{GC}}$ inflation references
overlaid. (\textbf{Bottom-right}) Three-tier recovery scorecard against
the Niu \emph{et al.}\ 2021 \cite{niu2021grain} 21-QTN catalogue
(Tier~1, blue), the Niu 2021 7 NEW candidate genes (Tier~2, green),
and the Ren \emph{et al.}\ 2023 \cite{ren2023rice} grain-quality
catalogue (Tier~3, red, normalised to the per-method maximum).
Extended caption with full annotation on the following page.}
\label{fig:sup_grain}
\end{figure}

\clearpage

\noindent\textbf{Supplementary Figure S7 (extended caption).}\hspace{1ex}%
The four-trait Manhattan tracks (Top panel) summarise the
PC1--PC10-corrected linear-model GWAS on the 3kRG panel
($N{=}1{,}847$ for TGW; $N{=}2{,}453$ for GL, GW, RLW; MAF~$\geq$~0.05).
The genome-wide significance threshold is $p \leq 5 \times 10^{-8}$
(red horizontal line); the suggestive threshold is
$p \leq 10^{-5}$ (orange). All four traits show the canonical
\emph{GS3}/\emph{qSW5}/\emph{GW5} grain-architecture peaks on Chr3 and
Chr5 plus several trait-specific QTNs (full breakdown in
Supplementary Table~S6).

\noindent\textbf{Reading the Q--Q plots (Bottom-left).}\hspace{1ex}%
Each point compares the $i$-th smallest observed p-value against the
expected uniform-null quantile ranking:
$x_i = -\log_{10}((i-0.5)/n)$,
$y_i = -\log_{10}(P_{(i)})$,
where $n \approx 27 \times 10^6$ is the number of variants per trait
after MAF~$\geq$~0.05 filtering. Three reading aids are overlaid on
every panel:
\emph{(i)} the dashed black diagonal $y = x$ is the expected position
of every point under a perfectly calibrated null (no association
anywhere in the genome); points on this line carry no signal.
\emph{(ii)} the dotted grey line at slope
$\sqrt{\lambda_{GC}}$ ($\lambda_{GC} \in [1.41, 1.78]$, displayed
per-panel) is where the \emph{bulk} of points should fall under
chi-square inflation alone, with no genuine association — the
classic genomic-control reference; the gap between $y=x$ and the
$\sqrt{\lambda_{GC}}$ line tells the reader how much of the upward
shift is attributable to inflation rather than signal.
\emph{(iii)} the trait-coloured cloud is sub-sampled by log-rank so
that the bulk and the upper tail both receive dense coverage at
every $-\log_{10}(\text{expected})$ value (the naive uniform-rank
sub-sample over-samples the tail and under-samples the bulk; see
Methods, ``\nameref{methods:multispecies}'').

\noindent\textbf{What the four Q--Q panels show.}\hspace{1ex}%
The bulk of points (lower-left) sits between the two reference lines,
indicating chi-square inflation $\lambda_{GC} > 1$ as expected for the
strongly stratified 3kRG population panel under PC1--PC10 correction
(XI/GJ subpopulation $F_{ST} \approx 0.5$). The upper-right tail
deflects sharply above both reference lines: variants that carry
true effects have much smaller p-values than even an inflated null
would predict, and their points climb to $-\log_{10}\!P$ values of
$\sim$150--250 for the strongest QTNs
(Chr3:16{,}733{,}441 / \emph{GS3} and
Chr5:5{,}371{,}529 / \emph{qSW5/GW5}, both visible at the upper-right
of every panel). Hundreds of points falling above the
$\sqrt{\lambda_{GC}}$ slope at expected $-\log_{10}(P) > 4$ indicate
hundreds of independent QTNs, consistent with the polygenic genetic
architecture of grain weight and shape (Niu \emph{et al.}\ 2021
\cite{niu2021grain} estimate narrow-sense heritability $h^2 = 0.88$
to $0.93$ for these traits). The fine-mapping pipeline applies
$\lambda_{GC}$ deflation ($z \to z / \sqrt{\lambda_{GC}}$) to the
per-trait z-scores before any credible-set computation, so the
inflated bulk visible in the Q--Q does not propagate into the
credible-set sizes reported in Supplementary Tables~S8 and S9.
Q--Q plots are rendered with rasterised markers; the underlying
p-value distributions are exact.

\noindent\textbf{Recovery scorecard (Bottom-right).}\hspace{1ex}%
Three-tier ground-truth recovery for each of the nine methods on the
combined 41-lead grain pass.
\emph{Tier~1} (blue, percentage scale): the Niu \emph{et al.}\ 2021
21-stable-QTN catalogue, scored as the QTN's chr:pos appearing in the
95\% credible set within $\pm$10\,kb of the QTN coordinate.
\emph{Tier~2} (green, percentage scale): the seven candidate genes
flagged as NEW by Niu 2021, scored at locus level (the candidate-gene
coordinate is covered by some method's credible set).
\emph{Tier~3} (red, normalised to the per-method maximum across the
nine methods): gene-level recovery against the Ren \emph{et al.}\ 2023
\cite{ren2023rice} grain-quality catalogue (104 grain-size +
grain-shape genes), with each method's count expressed as a percentage
of the maximum count observed for any method on this panel
(SBayesRC: 18 of $\sim$118 genes recovered = 100\%).
The strict variant-level numbers behind this panel are in
Supplementary Table~S8.
SBayesRC dominates the strict in-CS metric on Tier~1 (20/21 = 95.2\%,
on the strength of its full eigen-decomposed LD model); GAFM-MX,
HBP-MX and ENS dominate the strictest top-1-PIP exact-position
metric (10/21 = 47.6\%, the highest of any method tested,
exceeding SuSiE's 28.6\% and SBayesRC's 14.3\%).
Figure source: \texttt{tests/rice3k\_grain\_shape\_figures.py}.

\clearpage

\subsection*{Supplementary Table S5 --- IRRI 18-trait whole-genome scan summary}

\begin{table}[htbp]
\centering
\small
\caption{\textbf{IRRI 18-trait 3kRG whole-genome scan summary.}
$\lambda_{GC}$ is computed from $\sim$500\,000 random variants per
trait. ``GW indep.'' is the number of independent genome-wide-
significant leads after greedy 250-kb-window clustering on $p\leq
5\times10^{-8}$. ``CS$=$1 (GAFM)'' is the count among the 4
fine-mapped leads per trait for which GAFM narrows the 95\% credible
set to a single variant. The trait codes are IRRI Standard
Evaluation System ordinal scores; rough plain-language definitions
are given in the footnote.}
\label{tab:sup_irri_scan}
\begin{tabular}{lrrrrr}
\toprule
\textbf{Trait code} & \textbf{N} & $\lambda_{GC}$ & \textbf{min}~$p$ & \textbf{GW indep.} & \textbf{CS$=$1 (GAFM)} \\
\midrule
APCO\_REV\_REPRO  & 2261 & 2.07 & $3.1\times10^{-153}$ & 352  & 1 \\
AUCO\_REV\_VEG    & 2263 & 0.66 & $1.3\times10^{-244}$ & 1067 & 3 \\
AWCO\_REV         & 2264 & 0.95 & $1.9\times10^{-274}$ & 1067 & 2 \\
AWPR\_REPRO       & 2114 & 0.74 & $4.4\times10^{-90}$  & 1116 & 1 \\
BLCO\_REV\_VEG    & 2265 & 0.50 & $8.0\times10^{-55}$  & 1083 & 1 \\
CUAN\_REPRO       & 2264 & 1.62 & $2.6\times10^{-14}$  & 327  & 0 \\
CULT\_CODE\_REPRO & 2094 & 1.38 & $3.2\times10^{-34}$  & 1067 & 0 \\
CUST\_REPRO       & 2264 & 1.87 & $1.5\times10^{-24}$  & 559  & 0 \\
LA                & 1527 & 1.10 & $2.8\times10^{-11}$  & 81   & 0 \\
LIGCO\_REV\_VEG   & 2264 & 0.64 & $\sim$0              & 1031 & 2 \\
LLT\_CODE         & 2095 & 1.36 & $1.1\times10^{-24}$  & 882  & 0 \\
LPCO\_REV\_POST   & 2262 & 1.41 & $6.2\times10^{-35}$  & 517  & 0 \\
LPPUB             & 1526 & 1.38 & $1.3\times10^{-45}$  & 869  & 0 \\
LSEN              & 2263 & 1.37 & $1.5\times10^{-16}$  & 449  & 0 \\
PEX\_REPRO        & 2265 & 1.07 & $3.2\times10^{-14}$  & 168  & 0 \\
PTY               & 2262 & 0.93 & $4.4\times10^{-15}$  & 734  & 0 \\
SCCO\_REV         & 2115 & 1.81 & $5.1\times10^{-178}$ & 1054 & 1 \\
SPKF              & 2264 & 1.07 & $1.7\times10^{-10}$  & 74   & 0 \\
\midrule
\textbf{Total}    &      &      &                      &      & \textbf{11} \\
\bottomrule
\end{tabular}
\end{table}

\paragraph{Trait code definitions (IRRI Standard Evaluation System).}
APCO\_REV\_REPRO: apiculus colour at reproductive stage;
AUCO\_REV\_VEG: auricle colour at vegetative stage;
AWCO\_REV: awn colour;
AWPR\_REPRO: awn presence at reproductive stage;
BLCO\_REV\_VEG: blade (leaf) colour at vegetative stage;
CUAN\_REPRO: culm angle at reproductive stage;
CULT\_CODE\_REPRO: cultivar (varietal) code at reproductive stage;
CUST\_REPRO: culm strength at reproductive stage;
LA: leaf angle;
LIGCO\_REV\_VEG: ligule colour at vegetative stage;
LLT\_CODE: leaf-length-type code;
LPCO\_REV\_POST: leaf-pubescence colour post-reproductive;
LPPUB: leaf pubescence;
LSEN: leaf senescence;
PEX\_REPRO: panicle exsertion at reproductive stage;
PTY: panicle type;
SCCO\_REV: spikelet (grain pericarp) colour;
SPKF: spikelet fertility.

\clearpage

\subsection*{Supplementary Table S6 --- 3kRG grain weight + shape whole-genome scan summary (Niu \emph{et al.}\ 2021 phenotype reference panel)}

% Auto-generated by tests/rice3k_grain_shape_paper_tables.py
% 4-row addendum for Table S5 (IRRI 18-trait scan): grain weight + shape rerun
\begin{table}[htbp]
\centering\small
\caption{\textbf{3kRG grain weight + shape whole-genome scan summary (Niu \emph{et al.}\ 2021 phenotype reference panel).} Same plink2 \texttt{--glm linear} pipeline as Table~S5, with PC1--PC10 covariates from \texttt{rice\_3k.psam}. Niu \emph{et al.}\ 2021 \cite{niu2021grain} reports 21 stable QTNs from a compressed-MLM analysis on the same panel; we recover \textbf{20 of 21} (95\%) within 100\,kb of an independent lead at $p\leq5\times10^{-8}$. The fourth $\lambda_{GC}$ value is elevated above the 1.0 null ($\lambda_{GC} > 1.4$ for all four traits) under PC1--PC10 correction; this matches the IRRI scan and reflects the well-documented XI/GJ subpopulation structure ($F_{ST}\approx 0.5$).}
\label{tab:sup_grain_scan}
\begin{tabular}{lrrrr}
\toprule
\textbf{Trait} & \textbf{N indep.\ leads} & $\lambda_{GC}$ & \textbf{min}~$p$ & \textbf{Niu QTN match} \\
\midrule
TGW    &  1044 &  1.41 &        $8.6\times10^{-28}$ & 6/6 \\
GL    &  1102 &  1.78 &       $1.3\times10^{-193}$ & 4/5 \\
GW    &  1102 &  1.50 &       $6.2\times10^{-149}$ & 4/4 \\
RLW    &  1108 &  1.75 &       $4.8\times10^{-158}$ & 6/6 \\
\midrule
\textbf{Total} &       &       &  & \textbf{20/21} \\
\bottomrule
\end{tabular}
\end{table}

\paragraph{Phenotype provenance and quality control.}
TGW (thousand-grain weight, $N{=}1{,}847$) is from the
\texttt{thousand\_grain\_weight.txt} text file; GL, GW and RLW
($N{=}2{,}453$, $\mathrm{RLW}=\mathrm{GL/GW}$) are from the
\texttt{3K.shape.phe.xlsx} workbook --- both prepared from the
Niu \emph{et al.}\ 2021 \emph{BMC Genomics} \cite{niu2021grain}
multi-year, multi-environment GWAS panel (3 trial years; mean
across years used for the rerun). Sample identifiers
($\mathtt{B}/\mathtt{C}/\mathtt{IRIS\_313}$ prefixes) match the
3kRG plink2 \texttt{.psam} directly: shape ID coverage is 2453/2453
and TGW coverage is 1847/1847 against the 3{,}024-sample \texttt{.psam}
(joint coverage 1{,}832 for any pairwise correlation analysis).
Phenotype QC reproduces the four-trait pairwise Pearson
correlations of Niu Fig.~1b to within $\Delta r{<}0.04$ (e.g.\
$r(\mathrm{TGW},\mathrm{GW}){=}0.60$ vs.\ Niu's $0.63$;
$r(\mathrm{GL},\mathrm{RLW}){=}0.76$ vs.\ Niu's $0.76$;
$r(\mathrm{GW},\mathrm{RLW}){=}-0.78$ vs.\ Niu's $-0.78$). The
narrow-sense $h^2$ values reported by Niu are 0.92 (TGW), 0.88
(GL), 0.93 (GW), 0.92 (RLW) --- the highest in the rice
fine-mapping literature for QTN-supported trait classes.

\clearpage

\subsection*{Supplementary Table S7 --- Cross-species fine-mapping summary}

\begin{table}[htbp]
\centering
\footnotesize
\setlength{\tabcolsep}{4pt}
\caption{\textbf{Cross-species fine-mapping summary at baseline cache.}
GAFM CS $=$ 1 is the fraction of leads where flat-prior Bayesian GAFM
produces a single-variant 95\% credible set. ``HBP tighter'' is the
fraction of leads where HBP narrows the credible set further than GAFM
under the baseline cache (no per-variant continuous prior). Validation
uses a 250-kb window against species-specific literature catalogues
(Ren 2023 for rice; Bloom 2015 + Peter 2018 for yeast; AraGWAS plus
canonical flowering genes for Arabidopsis; GWAS-Catalog plus lipid-genetics
canon for human).}
\label{tab:sup_multispecies}
\begin{tabular}{@{}l@{\hspace{6pt}}r@{\hspace{6pt}}r@{\hspace{6pt}}r@{\hspace{4pt}}r@{\hspace{4pt}}r@{\hspace{4pt}}r@{\hspace{4pt}}r@{}}
\toprule
\textbf{Species} & \textbf{Traits} & \textbf{Loci}
                 & \shortstack[r]{\textbf{GAFM}\\\textbf{CS$=$1}}
                 & \shortstack[r]{\textbf{GAFM}\\\textbf{PIP$\geq 0.5$}}
                 & \shortstack[r]{\textbf{HBP}\\\textbf{tighter}}
                 & \shortstack[r]{\textbf{Both}\\\textbf{CS$=$1}}
                 & \textbf{Validated} \\
\midrule
Rice (3kRG, IRRI)      & 18 & 72  & 11 (15\%) & 16 (22\%) & 14 (19\%)  & 8 (11\%) & 24 (33\%) \\
Rice (3kRG, grain)$^\dagger$ &  4 & 41  & --        & --        & --         & --       & 20/21 (95.2\%) \\
Yeast (1011)           & 35 & 245 & 0 (0\%)   & 5 (2\%)   & 244 (100\%) & 0 (0\%) & 44 (18\%) \\
Arabidopsis (1001G)    & 14 & 54  & 3 (6\%)   & 7 (13\%)  & 29 (54\%)   & 2 (4\%)  & 3 (6\%)  \\
Human (Pan-UKB, GW)    & 3  & 321 & 8 (2\%)   & 47 (15\%) & 0 (0\%)     & 8 (2\%) & 21 (7\%) \\
\midrule
\textbf{Total}         & 74 & 733 & \textbf{23} & \textbf{79} & \textbf{315} & \textbf{19} & \textbf{112} \\
\bottomrule
\end{tabular}

\vspace{2pt}
\footnotesize{$^\dagger$ Validation count for the grain-shape pass uses the
panel-matched Niu \emph{et al.}\ 2021 \cite{niu2021grain} 21-stable-QTN
catalogue. The 20/21 (95.2\%) figure is panel-level (any GWAS lead at
$p\leq5\times10^{-8}$ within $\pm$100\,kb of the QTN). At the strict
variant-level (chr:pos in the 95\% credible set within $\pm$10\,kb),
the 9-method panel gives:
\textbf{SBayesRC 20/21 (95.2\%, 16 exact-position matches);
GAFM-MX = HBP-MX = ENS = SuSiE 12/21 (57.1\%);
HBP 10/21 (47.6\%);
GAFM, SuSiE-inf, FINEMAP-inf 8/21 (38.1\%) each}.
At the strictest top-1-PIP exact-position metric, GAFM-MX, HBP-MX,
and ENS each achieve \textbf{10/21 (47.6\%)} --- the highest of any
method, exceeding SuSiE (6/21, 28.6\%) and SBayesRC (3/21, 14.3\%). Per-method,
per-tier breakdown: Supplementary Table~S8. The mixture-prior step
adds a single SBayesRC-style 4-component Wakefield BF reweighting
on the LD-deconvolved, $\lambda_{\text{GC}}$-deflated z-scores
($\pi=(0.005, 0.003, 0.001, 0.001)$ within
non-zero components, $\gamma=(0.001, 0.01, 0.1, 1.0)$); the
deconvolution step is essential to preserve GAFM/HBP's LD-aware
ranking under the mixture, otherwise LD-correlated noise variants
with marginally larger marginal $|z|$ outcompete the causal during
multiplicative reweighting (Methods, Supplementary Table~S8\,c). The IRRI row's validation count uses Ren
\emph{et al.}\ 2023 \cite{ren2023rice} (gene-level recovery within
$\pm$250\,kb of any fine-mapped lead's top-PIP variant). Per-method,
per-tier recovery for the grain pass is in Supplementary Table~S8,
and the full nine-method per-locus credible-set comparison in
Supplementary Table~S9.
The grain pass spans 41 fine-mapped loci: top 5 GW-significant + top 2
suggestive per trait (28 discovery leads) plus 13 additional leads
matched within $\pm$100\,kb of a Niu 2021 QTN to ensure literature-anchored
coverage of the panel-matched catalogue.}
\end{table}

\clearpage

\subsection*{Supplementary Table S8 --- 3kRG grain weight + shape: nine-method ground-truth recovery scorecard}

% Auto-generated by tests/rice3k_grain_shape_paper_tables.py
\begin{table}[htbp]
\centering\small
\caption{\textbf{Three-tier ground-truth recovery on the 3kRG grain weight + shape rerun.} All six methods fine-map the same 41 leads (top 5 GW-significant + top 2 suggestive per trait, plus 13 Niu-augment leads; $\pm$100\,kb windows, 95\% credible sets). \emph{Tier 1:} variant-level recovery of the 21 stable QTNs reported in Niu \emph{et al.}\ 2021 \cite{niu2021grain}, scored as the QTN's chr:pos appearing in the 95\% credible set of any locus assigned to the matching trait. \emph{Tier 2:} locus-level recovery of the 7 NEW candidate genes flagged by Niu \emph{et al.}\ 2021 (anchor: candidate-gene QTN coordinate covered by the credible set). \emph{Tier 3:} gene-level recovery against the Ren \emph{et al.}\ 2023 \cite{ren2023rice} grain-quality catalogue (104 grain-size + grain-shape genes), scored as the gene falling within $\pm$250\,kb of the top-PIP variant of any matching-trait locus.}
\label{tab:sup_grain_recovery}
\begin{tabular}{lrrrr}
\toprule
\textbf{Method} & \shortstack[r]{\textbf{Tier 1}\\\textbf{Niu QTN}\\\textbf{(in CS)}} & \shortstack[r]{\textbf{Tier 1}\\\textbf{Niu QTN}\\\textbf{top-1 PIP}} & \shortstack[r]{\textbf{Tier 2}\\\textbf{NEW gene}\\\textbf{(locus)}} & \shortstack[r]{\textbf{Tier 3}\\\textbf{Ren genes}\\\textbf{(±250\,kb)}} \\
\midrule
GAFM           &    8/21 (38\%) &       0/21 &     3/7 (43\%) &    17 \\
HBP            &   10/21 (48\%) &       0/21 &     4/7 (57\%) &    17 \\
GAFM-MX        &   12/21 (57\%) &      10/21 &     6/7 (86\%) &    18 \\
HBP-MX         &   12/21 (57\%) &      10/21 &     6/7 (86\%) &    18 \\
ENS            &   12/21 (57\%) &      10/21 &     6/7 (86\%) &    18 \\
SuSiE          &   12/21 (57\%) &       6/21 &     6/7 (86\%) &    17 \\
SuSiE-inf      &    8/21 (38\%) &       3/21 &     6/7 (86\%) &    17 \\
FINEMAP-inf    &    8/21 (38\%) &       3/21 &     5/7 (71\%) &    16 \\
SBayesRC       &   20/21 (95\%) &       3/21 &     6/7 (86\%) &    18 \\
\bottomrule
\end{tabular}
\end{table}

\paragraph{Method panel.} The five methods are GAFM (graph-augmented
fine-mapping; this work, $\alpha{=}0.7$, sumstats interface,
\texttt{l1\_finemap\_from\_sumstats}); HBP (hierarchical belief
propagation; this work, sumstats interface,
\texttt{hbp\_finemap\_from\_sumstats}); SuSiE
\cite{wang2020simple} via \texttt{susieR::susie\_rss} on signed
in-sample LD; SuSiE-inf \cite{cui2024improving} via the
FinucaneLab Python \texttt{susieinf} reference implementation; and
FINEMAP-inf \cite{cui2024improving} via the FinucaneLab Python
\texttt{finemapinf} reference implementation. All five share the
same 250\,kb LD windows and 95\% credible-set coverage; GAFM/HBP
additionally use the rice multi-omics graph cache (RAP-DB
variant--gene + Ren 2023 gene--pathway + RicePPINet PPI). The
SuSiE-inf and FINEMAP-inf baselines are the state-of-the-art
infinitesimal-effects methods from Wu \emph{et al.}\ 2026 (Fig.~4b).
SBayesRC \cite{zheng2024leveraging} is included as the 6th method.
We built a 3kRG-specific eigen-decomposed LD reference from scratch
using GCTB\cite{zeng2018gctb}~2.05beta + the R \texttt{LDstep1}--\texttt{LDstep4}
pipeline over 23 non-overlapping blocks covering all 41 fine-mapped
leads (188\,K MAF~$\geq$~0.05 variants; ~10 min wall-clock).
SBayesRC achieves \textbf{20/21 (95.2\%)} variant-level recovery
of the Niu QTNs in 95\% credible sets within $\pm$10\,kb (16 exact
position matches), the strongest of all six methods on this rice
grain pass.

\clearpage

\subsection*{Supplementary Table S9 --- 3kRG grain pass: full nine-method per-locus credible-set comparison}

% Auto-generated by tests/rice3k_grain_shape_paper_tables.py
\begin{table}[htbp]
\centering\footnotesize
\setlength{\tabcolsep}{4pt}
\caption{\textbf{Per-locus credible-set comparison on the 3kRG grain weight + shape rerun.} Cells report 95\% credible-set size / top-variant PIP. We display the curated 5-method best-of-class subset; the full 9-method comparison (GAFM, HBP, GAFM-MX, HBP-MX, ENS, SuSiE, SuSiE-inf, FINEMAP-inf, SBayesRC) is in \texttt{data/rice\_3k/results/grain\_finemap/grain\_finemap\_summary.tsv}. All methods receive $\lambda_{GC}$-deflated z-scores (z $\to$ z/$\sqrt{\lambda_{GC}}$). Loci selected as the top 5 GW-significant + top 2 suggestive per trait, plus 13 leads matched within $\pm$100\,kb of a Niu 2021 QTN. $N_v$ is the number of variants in the $\pm$100\,kb fine-map window.}
\label{tab:sup_grain_per_locus}
\begin{tabular}{l l r r r r r r r}
\toprule
\textbf{Trait} & \textbf{Locus} & \textbf{lead $p$} & $N_v$ & \textbf{GAFM-MX} & \textbf{HBP-MX} & \textbf{ENS} & \textbf{SuSiE} & \textbf{SBayesRC} \\
\midrule
\addlinespace[0.5ex]
GL & Chr3:16,733,441 & $1.3{\times}10^{-193}$ & 4395 & 1/1.00 & 1/1.00 & 1/1.00 & 1/1.00 & 4027/1.00 \\
GL & Chr3:17,013,703 & $8.7{\times}10^{-80}$ & 2959 & 1/1.00 & 1/1.00 & 1/1.00 & 1/1.00 & 2708/1.00 \\
GL & Chr6:14,837,641 & $1.3{\times}10^{-72}$ & 4782 & 2/0.65 & 2/0.66 & 2/0.66 & 1/1.00 & 4337/1.00 \\
GL & Chr1:38,470,016 & $6.0{\times}10^{-67}$ & 2090 & 1/0.97 & 1/0.98 & 1/0.98 & 2/1.00 & 1901/1.00 \\
GL & Chr6:21,149,455 & $1.4{\times}10^{-66}$ & 2840 & -- & -- & -- & -- & 2589/0.49 \\
GL & Chr6:21,149,455 & $1.4{\times}10^{-66}$ & 2840 & 2/0.93 & 2/0.94 & 2/0.94 & 1/1.00 & -- \\
GL & Chr5:5,371,949 & $4.3{\times}10^{-45}$ & 2858 & 13/0.42 & 13/0.43 & 13/0.43 & 1/1.00 & 2618/0.66 \\
GL & Chr7:28,291,625 & $1.0{\times}10^{-28}$ & 1148 & 6/0.35 & 6/0.36 & 6/0.35 & 1/1.00 & 1050/0.25 \\
GL & Chr4:29,310,722 & $1.8{\times}10^{-19}$ & 2484 & 5/0.77 & 5/0.79 & 5/0.78 & 1/1.00 & 2263/0.98 \\
GL & Chr5:28,778,827 & $5.3{\times}10^{-8}$ & 1796 & 581/0.25 & 131/0.31 & 254/0.28 & 2/1.00 & 1654/0.22 \\
GL & Chr12:26,216,487 & $6.7{\times}10^{-8}$ & 2972 & 2321/0.31 & 1499/0.42 & 2069/0.36 & 1/1.00 & 2712/1.00 \\
\addlinespace[0.5ex]
GW & Chr5:5,371,686 & $6.2{\times}10^{-149}$ & 2851 & 1/0.98 & 1/0.98 & 1/0.98 & 1/1.00 & 2648/1.00 \\
GW & Chr5:5,007,414 & $8.0{\times}10^{-59}$ & 1878 & 1/1.00 & 1/1.00 & 1/1.00 & 1/1.00 & 1739/1.00 \\
GW & Chr7:24,943,451 & $3.5{\times}10^{-56}$ & 1729 & 920/0.39 & 250/0.49 & 630/0.44 & 1/1.00 & 1602/0.97 \\
GW & Chr7:25,214,651 & $2.5{\times}10^{-54}$ & 3882 & 2/0.88 & 2/0.88 & 2/0.88 & 1/1.00 & 3600/0.98 \\
GW & Chr7:24,629,753 & $3.8{\times}10^{-54}$ & 3158 & 1/1.00 & 1/1.00 & 1/1.00 & 1/1.00 & 2937/0.82 \\
GW & Chr3:16,733,441 & $1.3{\times}10^{-25}$ & 4395 & -- & -- & -- & -- & 4077/0.90 \\
GW & Chr3:16,733,441 & $1.3{\times}10^{-25}$ & 4395 & 4/0.58 & 4/0.59 & 4/0.59 & 1/1.00 & -- \\
GW & Chr8:26,530,724 & $8.6{\times}10^{-18}$ & 2440 & 1/0.98 & 1/0.99 & 1/0.98 & 1/1.00 & 2262/1.00 \\
GW & Chr4:28,798,865 & $1.1{\times}10^{-7}$ & 1954 & 4/0.67 & 4/0.69 & 4/0.68 & 1/1.00 & 1805/1.00 \\
GW & Chr11:5,112,561 & $2.2{\times}10^{-7}$ & 4870 & 3394/0.15 & 1561/0.20 & 2818/0.17 & 1/1.00 & 4515/0.91 \\
\addlinespace[0.5ex]
RLW & Chr3:16,733,441 & $4.8{\times}10^{-158}$ & 4395 & 1/1.00 & 1/1.00 & 1/1.00 & 1/1.00 & 3873/0.96 \\
RLW & Chr5:5,371,609 & $1.6{\times}10^{-146}$ & 2850 & 3/0.61 & 3/0.62 & 3/0.61 & 1/1.00 & 2533/1.00 \\
RLW & Chr3:17,006,465 & $3.9{\times}10^{-61}$ & 3187 & 1/1.00 & 1/1.00 & 1/1.00 & 1/1.00 & 2778/0.99 \\
RLW & Chr6:14,837,641 & $1.9{\times}10^{-58}$ & 4782 & 3/0.91 & 2/0.94 & 2/0.92 & 1/1.00 & 4106/1.00 \\
RLW & Chr7:24,902,815 & $6.9{\times}10^{-58}$ & 2075 & 1/0.99 & 1/0.99 & 1/0.99 & 1/1.00 & 1805/0.99 \\
RLW & Chr7:24,629,753 & $6.5{\times}10^{-57}$ & 3158 & 28/0.79 & 21/0.82 & 24/0.81 & 1/1.00 & 2790/0.99 \\
RLW & Chr8:26,552,429 & $3.8{\times}10^{-23}$ & 2401 & -- & -- & -- & -- & 2110/1.00 \\
RLW & Chr8:26,552,429 & $3.8{\times}10^{-23}$ & 2401 & 1/0.96 & 1/0.96 & 1/0.96 & 1/1.00 & -- \\
RLW & Chr4:29,310,722 & $3.3{\times}10^{-21}$ & 2484 & 5/0.57 & 5/0.59 & 5/0.58 & 1/1.00 & 2173/1.00 \\
RLW & Chr1:3,691,963 & $1.6{\times}10^{-15}$ & 2412 & 501/0.23 & 313/0.26 & 432/0.25 & 1/1.00 & 2128/1.00 \\
RLW & Chr2:3,273,658 & $2.7{\times}10^{-7}$ & 1966 & 614/0.48 & 143/0.57 & 346/0.52 & 1/1.00 & 1755/0.94 \\
RLW & Chr1:40,718,057 & $2.9{\times}10^{-7}$ & 2136 & 1913/0.25 & 1632/0.37 & 1827/0.31 & 1/1.00 & 1928/0.22 \\
\addlinespace[0.5ex]
TGW & Chr3:16,733,441 & $8.6{\times}10^{-28}$ & 4395 & 1/1.00 & 1/1.00 & 1/1.00 & 1/1.00 & 4090/0.99 \\
TGW & Chr7:28,323,091 & $1.9{\times}10^{-27}$ & 1212 & 5/0.63 & 5/0.65 & 5/0.64 & 1/1.00 & 1124/0.51 \\
TGW & Chr9:10,266,788 & $2.6{\times}10^{-25}$ & 5952 & 7/0.31 & 7/0.31 & 7/0.31 & 1/1.00 & 5520/0.89 \\
TGW & Chr9:20,921,346 & $5.4{\times}10^{-25}$ & 2043 & 2/0.54 & 2/0.55 & 2/0.54 & 1/1.00 & 1904/0.20 \\
TGW & Chr3:18,966,681 & $1.2{\times}10^{-24}$ & 3874 & 6/0.29 & 6/0.30 & 6/0.29 & 1/1.00 & -- \\
TGW & Chr3:18,966,681 & $1.2{\times}10^{-24}$ & 3874 & -- & -- & -- & -- & 3601/0.65 \\
TGW & Chr9:21,335,551 & $2.6{\times}10^{-24}$ & 2350 & 3/0.52 & 3/0.52 & 3/0.52 & 1/1.00 & 2178/1.00 \\
TGW & Chr3:4,431,883 & $8.4{\times}10^{-23}$ & 2524 & 47/0.26 & 43/0.29 & 45/0.27 & 1/1.00 & 2342/1.00 \\
TGW & Chr5:5,299,051 & $4.1{\times}10^{-18}$ & 2266 & 17/0.66 & 15/0.69 & 16/0.68 & 1/1.00 & 2102/1.00 \\
TGW & Chr11:3,051,432 & $1.4{\times}10^{-14}$ & 2262 & 56/0.58 & 6/0.62 & 25/0.60 & 65/1.00 & 2094/0.77 \\
TGW & Chr2:30,410,561 & $5.0{\times}10^{-8}$ & 2410 & 1180/0.57 & 404/0.72 & 855/0.65 & 1/1.00 & 2249/0.10 \\
TGW & Chr11:2,226,393 & $6.0{\times}10^{-8}$ & 2211 & 5/0.24 & 5/0.24 & 5/0.24 & 1/1.00 & 2043/1.00 \\
\bottomrule
\end{tabular}
\end{table}

\clearpage

\subsection*{Supplementary Table S10 --- Cross-species CS=1 sharpening of GAFM-MX, HBP-MX and ENS on real GWAS leads}

\begin{table}[htbp]
\centering\footnotesize
\setlength{\tabcolsep}{4pt}
\caption{\textbf{Cross-species CS$=$1 sharpening of GAFM-MX, HBP-MX and ENS on real GWAS leads.}
For each species, we run all five methods (GAFM, HBP, GAFM-MX, HBP-MX, ENS)
on the species' actual lead-locus catalogue (yeast: 245 leads from 35
1011-Genomes traits with grammar-corrected sumstats; Arabidopsis:
54 leads from 14 AraPheno traits; IRRI rice: 72 leads from 18 3kRG
traits) using the same multi-omics graph cache
as the rice 3kRG pass. ``CS$=$1'' is the fraction of leads where the 95\%
credible set is a single variant; ``median top PIP'' is the median PIP
of the top-PIP variant across leads. The mixture-prior step (MX/ENS)
delivers $20$--$60\times$ sharpening of the credible set vs base GAFM/HBP
on real leads, exceeding the chr22 single-causal-F1 simulation ratio
(Supplementary Table~S11) because real GWAS lead loci have stronger marginal
signals than the simulated single-causal panel and the mixture BF is most
discriminative at large $|z|$. Reproduce with
\texttt{python tests/yeast\_finemap\_all.py},
\texttt{CACHE\_VERSION=v2 python tests/arabidopsis\_finemap.py}, and
\texttt{python tests/rice3k\_irri\_finemap.py}.}
\label{tab:sup_v15_cross_species}
\begin{tabular}{l@{\hspace{6pt}}r@{\hspace{4pt}}r@{\hspace{6pt}}r@{\hspace{4pt}}r@{\hspace{6pt}}r@{\hspace{4pt}}r}
\toprule
                 & \multicolumn{2}{c}{\textbf{Yeast (1011)}}
                 & \multicolumn{2}{c}{\textbf{Arabidopsis (1001G)}}
                 & \multicolumn{2}{c}{\textbf{IRRI rice (3kRG)}} \\
                 & \multicolumn{2}{c}{$n=245$ leads}
                 & \multicolumn{2}{c}{$n=54$ leads}
                 & \multicolumn{2}{c}{$n=72$ leads} \\
\cmidrule(lr){2-3}\cmidrule(lr){4-5}\cmidrule(lr){6-7}
\textbf{Method}  & CS$=$1 & med top PIP & CS$=$1 & med top PIP & CS$=$1 & med top PIP \\
\midrule
GAFM             &   0/245 (0\%)    & 0.009 &   2/54 (3.7\%)   & 0.041 &  11/72 (15.3\%) & 0.005 \\
HBP              &   0/245 (0\%)    & 0.014 &   2/54 (3.7\%)   & 0.106 &   8/72 (11.1\%) & 0.018 \\
GAFM-MX          &  60/245 (24.5\%) & 0.356 &  35/54 (64.8\%)  & 0.996 &  30/72 (41.7\%) & 0.837 \\
HBP-MX           &  63/245 (25.7\%) & 0.412 &  35/54 (64.8\%)  & 0.997 &  30/72 (41.7\%) & 0.855 \\
ENS              &  60/245 (24.5\%) & 0.387 &  35/54 (64.8\%)  & 0.996 &  30/72 (41.7\%) & 0.846 \\
\bottomrule
\end{tabular}
\end{table}

\clearpage

\subsection*{Supplementary Table S11 --- Chr22 head-to-head across signal strength (eight-method panel)}

\begin{table}[htbp]
\centering\small
\caption{\textbf{Chr22 single-causal F1 head-to-head across $h^2$ (eight-method panel).}
$N{=}30$ replicate windows (50\,kb, MAF~$\geq$~0.05 in 1000\,Genomes
Phase~3 chr22) at each of three heritability levels. ``rank-1'' is the
fraction of replicates where the causal variant is the
single highest-PIP variant; ``mean PIP'' is averaged at the causal
variant. The mixture-prior methods (GAFM-MX, HBP-MX, ENS) achieve
the same rank-1 rate as base GAFM/HBP but deliver
$2$--$3\times$ the confidence at the causal variant; sharpening is
largest at weakest signal (where fine-mapping needs help most). All
GAFM/HBP variants run in $\leq 5$\,ms per locus, $\geq 100\times$
faster than SuSiE/SuSiE-inf/FINEMAP-inf. Reproduce with
\texttt{N\_REPS=30 H2=\{0.02,0.05,0.10\} python tests/benchmark\_v15\_chr22.py}.}
\label{tab:sup_v15_chr22}
\begin{tabular}{l c c c c c c}
\toprule
 & \multicolumn{2}{c}{\textbf{$h^2 = 0.02$}}
 & \multicolumn{2}{c}{\textbf{$h^2 = 0.05$}}
 & \multicolumn{2}{c}{\textbf{$h^2 = 0.10$}} \\
\cmidrule(lr){2-3}\cmidrule(lr){4-5}\cmidrule(lr){6-7}
\textbf{Method} & rank-1 & mean PIP & rank-1 & mean PIP & rank-1 & mean PIP \\
\midrule
GAFM         & 14/30 & 0.153 & 16/30 & 0.324 & 17/30 & 0.416 \\
HBP          & 14/30 & 0.252 & 16/30 & 0.398 & 17/30 & 0.464 \\
GAFM-MX      & 14/30 & \textbf{0.435} & 16/30 & \textbf{0.534} & 17/30 & \textbf{0.578} \\
HBP-MX       & 14/30 & \textbf{0.436} & 16/30 & \textbf{0.534} & 17/30 & \textbf{0.578} \\
ENS          & 14/30 & \textbf{0.436} & 16/30 & \textbf{0.534} & 17/30 & \textbf{0.578} \\
SuSiE        & 18/30 & 0.549 & 21/30 & 0.674 & 22/30 & 0.747 \\
SuSiE-inf    & 18/30 & 0.543 & 22/30 & 0.693 & 24/30 & 0.765 \\
FINEMAP-inf  & 19/30 & 0.556 & 21/30 & 0.699 & 22/30 & 0.788 \\
\bottomrule
\end{tabular}
\end{table}
